% Overleaf: XeLaTeX, главный файл report.tex.
\documentclass[11pt,a4paper]{article}
\usepackage[margin=23mm]{geometry}
\usepackage{fontspec}
\IfFontExistsTF{FreeSerif}{\setmainfont{FreeSerif}}{\setmainfont{Times New Roman}}
\usepackage[english,russian]{babel}
\usepackage{array,tabularx}
\setlength{\parindent}{0pt}
\setlength{\parskip}{5pt}
\renewcommand{\arraystretch}{1.35}

\newcommand{\partofreport}[3]{\clearpage\section{#1}\textit{Ответственный: #2. Лимит: #3}\par}
\newcommand{\slot}[1]{\par\noindent\fbox{\parbox[c][13mm][c]{.96\linewidth}{\centering #1}}\par}

\begin{document}

{\Large\bfseries Моделирование перекрёстка\par}
\section{Сводка результатов}
\textit{Ответственный: Егор Турыгин (R-3). Лимит: 1 стр.}

\textbf{Таблица 6 — Сводка результатов}\par
\begin{tabularx}{\linewidth}{|p{.32\linewidth}|X|}
\hline
\textbf{Строка} & \textbf{Результат} \\ \hline
Вариант & \\ \hline
Методика & \\ \hline
Базовый план & \\ \hline
Выбранный план & \\ \hline
Формула Уэбстера & \\ \hline
Задание 5 & \\ \hline
Задание 6 & \\ \hline
Бонусы & \\ \hline
\end{tabularx}

% Вариант: команда, спрос, пиковый въезд. Методика: W, T, число прогонов и SEED.
% Базовый план: задержка с интервалом и максимальная очередь.
% Выбранный план: CYCLE, NS_GREEN, EW_GREEN и задержка с интервалом на holdout.

\partofreport{Задание 1. Независимые случайные потоки}{Владислав Полянский (T-11), Георгий Дмитриев (T-12), Мира Кру (T-13), Артур Лосин (T-14)}{1 стр.}
Правило вывода SEED и результат теста идентичности журналов.
\slot{Фрагмент журнала из пяти строк}
\slot{Стандартные отклонения 20 парных разностей с общими и независимыми потоками}
Когда общие потоки уменьшают разброс и когда эффект может отсутствовать.

\partofreport{Задание 2. Показатели качества}{Егор Турыгин (T-21), Мира Кру (T-22), Сергей Новиков (T-24)}{2 стр.}
Определение ожидания, формула накопленной задержки и выборка для 95-го процентиля.
\slot{Показатели базового плана по въездам}
\slot{Сверка с ручным подсчётом и суммы по въездам}

\partofreport{Задание 3. Точность результатов}{Артур Лосин (T-32), Мира Кру (T-33), Сергей Новиков (T-34)}{2 стр.}
\slot{График прогрева и скользящего среднего с отметкой W}
\slot{Среднее, стандартное отклонение и 95\% ДИ для базового плана и парной разности}
Расчёт и проверка необходимого числа прогонов.

\partofreport{Задание 4. Подбор плана светофора}{Егор Турыгин (T-41), Георгий Дмитриев (T-42), Мира Кру (T-44), Владислав Полянский (T-45), Сергей Новиков (T-46)}{3 стр.}
Целевая функция, ограничения и правило выбора. Итоги перебора и ускоренного поиска.
\slot{Тепловая карта по длине цикла и доле зелёного}
\slot{Отложенная проверка базового и трёх лучших планов с интервалами}
Расчёт цикла Уэбстера и сравнение с выбранным циклом.

\partofreport{Задание 5. Новая схема фаз}{Артур Лосин (T-51), Георгий Дмитриев (T-52), Егор Турыгин (T-53)}{1 стр.}
\slot{Диаграмма или таблица фаз и результат проверки сигналов}
\slot{Сравнение задержек, очередей и конфликтов с двухфазным планом с интервалами}
Вывод в двух-трёх предложениях.

\partofreport{Задание 6. Пиковый спрос и адаптивное управление}{Георгий Дмитриев (T-61), Сергей Новиков (T-62), Егор Турыгин (T-63), Мира Кру (T-64)}{2 стр.}
\slot{Очереди по въездам для обеих политик с отмеченным окном пика}
\slot{Задержка, максимальная очередь и время восстановления для обеих политик}
Данные контроллера, правило продления, минимум и максимум зелёного.

\partofreport{Бонусные задания}{Мира Кру (B-1), Артур Лосин (B-2), Сергей Новиков (B-3)}{1 стр.}
\slot{Отложенная проверка для каждого выполненного бонуса}
Указать, изменились ли выводы. Убрать раздел, если бонусов нет.

\partofreport{Выводы и ограничения модели}{Егор Турыгин (R-3)}{1 стр.}
Обоснование выбранного плана, существенные ограничения и область применимости результатов.

\partofreport{Вклад участников}{Сергей Новиков (R-6)}{1 стр.}
\begin{tabularx}{\linewidth}{|p{.27\linewidth}|p{.24\linewidth}|X|}
\hline
\textbf{Участник} & \textbf{Роль} & \textbf{Что сделал} \\ \hline
 & & \\ \hline
 & & \\ \hline
 & & \\ \hline
 & & \\ \hline
 & & \\ \hline
 & & \\ \hline
\end{tabularx}

\end{document}
